% Reconstructed after the workspace reset from the audited final text and retained conversation.
% Root certified the mathematics at 09:58:11.544224 UTC; Frobenius independently passed.
\paragraph{C.3an (Refuted, with an explicit integral counterexample): an infinite top quotient in characteristic five.}
Let $k$ be any algebraically closed field of characteristic $5$.
We specify an excellent normal strictly henselian local $k$-algebra
$R$, with residue field $k$ and dimension three, whose puncture $X$
is regular, and construct a surjection
\[
 B^2(X):=\mathrm{CH}^2(X)/\mathrm{CH}^2(X)_{\mathrm{alg}}
                          \longrightarrow(k,+).
\]
Here algebraic equivalence uses actual families on fixed products.
Consequently this integral top quotient is infinite, including over
$\overline{\mathbb F}_5$. The exhibited additive quotient has exponent
five. No assertion about the rational or complex top quotient follows.

\begin{proof}
\emph{1. The surface and its ample cotangent line.}
In homogeneous coordinates $(Y:Z:W)$ put
\[
 C:\quad Y^5-YW^4=Z^4W,\qquad
 \infty=(0:1:0),\qquad D_1=2\infty.
\]
The affine $Y$ derivative is $-1$; at the unique point at infinity
the $W$ derivative is $-Z^4$. Thus $C$ is smooth of genus six.
On $W\ne0$ the equation is $y^5-y=z^4$, so $z$ is an etale
coordinate and $dz$ is nowhere zero. Its divisor is
$10\infty=5D_1$ by the degree of the canonical divisor.
Use the Frobenius sequence
\[
 0\longrightarrow\mathcal O_C(-D_1)
 \longrightarrow F_*\mathcal O_C(-5D_1)
 \longrightarrow B_C(-D_1)\longrightarrow0,
 \qquad B_C=F_*\mathcal O_C/\mathcal O_C.
\]
The normalized exact differential $dz$ defines a nonzero section
of the last sheaf. Its connecting class $\alpha$ is nonzero,
because the two negative line bundles have no global sections,
and is killed by Frobenius. Let
\[
 0\longrightarrow\mathcal O_C(-D_1)\longrightarrow E
                 \longrightarrow\mathcal O_C\longrightarrow0
                                                        \tag{C.3an.1}
\]
be its extension. This is the explicit Tango construction of
\citet[Example~1.3, Lemma~1.4 and Section~1.3, author pp.~3--6]
{D_Mukai2011}.

Let $P=\mathbb P_C(E)$ be the bundle of lines. Its section
$S_\infty$ from the subline has normal bundle
$\mathcal O_C(D_1)$. The Frobenius pullback has a unique
splitting, since $H^0(C,\mathcal O_C(-5D_1))=0$.
The inverse image of its complementary section is the smooth
integral divisor $G$, disjoint from $S_\infty$, finite purely
inseparable of degree five over $C$. On the affine chart its
equation is $s^5=z$; the normalized nonvanishing differential
verifies smoothness on every chart. The element $z$ is not
a fifth power in $k(C)$ because $dz\ne0$.

Write $\bar g:P\to C$. Fibre degree and restriction to the
section give, with its sign determined rather than assumed,
\[
                G\sim5S_\infty-\bar g^*(10\infty).
                                                        \tag{C.3an.2}
\]
Indeed $G-5S_\infty$ is pulled back from $C$ and restricts to
$\mathcal O_C(-5D_1)$. Set
\[
 M=\mathcal O_P(-3S_\infty)\otimes\bar g^*\mathcal O_C(5\infty),
 \qquad V=\operatorname{Spec}_P(\mathcal O_P\oplus M),
                                                        \tag{C.3an.3}
\]
where $M^2\to\mathcal O_P$ has divisor $G+S_\infty$ and scalar
fixed by $T^2=s^5-z$. The cover is connected by its simple
branch valuation, smooth because its two branch curves are
smooth and disjoint and $2$ is invertible, and projective.
Let $g:V\to C$, let $S$ be the reduced inverse image of
$S_\infty$, and define
\[
                 D=S+g^*\infty,\qquad L=\mathcal O_V(5D).
                                                        \tag{C.3an.4}
\]
The normal bundle $N_{S/V}=M^{-1}|_{S_\infty}$ is
$\mathcal O_C(\infty)$, so $S^2=1$ and $D^2=3$.
Every fibre is the integral compactification of
$T^2=(s-c)^5$, meeting $S$ once. Hence $D\cdot S=2$ and
$D\cdot F=1$ for a fibre. Every other horizontal integral
curve has nonnegative intersection with $S$ and positive
degree over $C$. Nakai--Moishezon proves $D$ and $L$ ample;
compare \citet[Lemma~2.2, author p.~8]{D_Mukai2011}.

The differential $dT$ gives a nowhere vanishing section
\[
               \eta\in H^0(V,\Omega^1_{V/k}\otimes L^{-1}).
                                                        \tag{C.3an.5}
\]
Here are the exact orders. Near infinity put $q=Y/Z$ and
$b=W/Z$. The equation is $b=q^5-qb^4$, so $q$ is a
uniformizer and $b$ has order five. Direct subtraction gives
\[
             z-q^{-5}=b^3/q^4=q^{10}c,\qquad c=b^3/q^{14}.
\]
The function $c$ is regular of order one and $dc$ is a
generator. The normalized coordinates
$s=q^2s_i+q^{-1}$ and $T=q^5T_i$ satisfy
$T_i^2=s_i^5-c$ and $dT=q^5dT_i$.
More generally these same formulas hold on normalized
charts with a local equation for $\infty$ in place of $q$.
On a finite fibre chart $s_i,T_i$ are smooth coordinates.
At the infinity section, $x=1/s_i$ and $w=T_i/s_i^3$ give
\[
                 w^2=x-cx^6,\qquad dT_i=-\frac{x^3}{2w}\,dc.
\]
Since $x/w^2$ is a unit, the order along $S$ is five.
These identities give the exact divisor $5S+5g^*\infty$,
with no further zero at the intersection, proving (C.3an.5).
They also independently verify the cotangent assertion of
\citet[Proposition~2.3 and equation~(25), author pp.~8,13]
{D_Mukai2011}.

\emph{2. The actual cone and resolution.}
Take the full ample section ring and its homogeneous maximal ideal:
\[
 A=\bigoplus_{n\ge0}H^0(V,L^n),\quad \mathfrak n=A_{>0},\quad
 R=(A_{\mathfrak n})^{\mathrm{sh}},\quad
 X=\operatorname{Spec}R\setminus\{\mathfrak m_R\}.
                                                        \tag{C.3an.6}
\]
The geometric point has the specified residue field $k$.
The section ring is finitely generated and normal: its
embedding in $k(V)[u]$ identifies it with the intersection
of the coefficient valuation conditions for regular sections.
It has dimension three. The morphism
$Z_0=\operatorname{Tot}(L^{-1})\to\operatorname{Spec}A$
is proper. Indeed choose a sufficiently large Veronese
whose ample line bundle is projectively normal and has
degree-one-generated section ring. The fibre-power map
to its usual cone resolution is finite, and that resolution
is proper. The resulting properness over the Veronese
cone implies properness over $\operatorname{Spec}A$, which
is finite and separated over the Veronese cone.
Off the zero section, consecutive sufficiently large
powers of $L$ provide an invertible degree-one ratio.
Their degree-zero localizations identify the morphism
with the usual $\mathbb G_m$-bundle. Thus it is an
isomorphism away from the vertex and contracts precisely $V$.
No generation assertion about $L$ itself is required.

Put $Z=Z_0\times_A\operatorname{Spec}R$. It is regular,
proper over $R$, with complement $X$ of its exceptional
divisor $V$, and retains the projection $\pi:Z\to V$.
Its conormal is $L$, and the first neighborhood is split:
\[
          V_2=\operatorname{Spec}_V(\mathcal O_V\oplus L),
                            \qquad L^2=0.
                                                        \tag{C.3an.7}
\]
Strict henselization does not alter this neighborhood.
Its map to the cone factors through an Artinian quotient
supported at the vertex, already strictly henselian with
residue field $k$. This also applies to every finite
neighborhood; no algebraization of arbitrary formal classes
is asserted.

\emph{3. The integral cohomological invariant.}
The exact calculation of
\citet[unnumbered theorem and corollary, p.~A1204]{frVanDerKallen1971},
equivalently \citet[Chapter~III, Exercise~5.14(a)--(c), pp.~46--47]
{frWeibelKbookIII}, gives the split sheaf decomposition
\[
 \mathcal K_{2,V_2}=\mathcal K_{2,V}
                       \oplus(\Omega^1_{V/k}\otimes L).
                                                        \tag{C.3an.8}
\]
This is integral, requiring only $1/2$. We choose the
negative of the Exercise~5.14 map in Weibel's modern
Dennis--Stein convention, so $\{1+tc,u\}$ maps to
$tc\,d\log u$. The original differential module is absolute
over $\mathbb Z$. Since every scalar of $k$ is a fifth
power, $\Omega^1_{k/\mathbb Z}=0$ and that module is exactly
$\Omega^1_{V/k}$. Rescaling the square-zero parameter
gives the displayed line-bundle twist.

Quillen's integral Bloch formula
\citep[Section~7, Theorems~5.11 and~5.19, pp.~133--137]
{rootQuillen1973} gives
$\mathrm{CH}^2(Z)=H^2(Z,\mathcal K_2)$.
It applies by passage from the smooth finite-type pointed
etale models defining $R$; finite cycle, rational-equivalence
and Gersten-complex data commute with this limit.
Restriction to $V_2$ and the split projection define
\[
 \delta:\mathrm{CH}^2(Z)\longrightarrow
                              H^2(V,\Omega^1_{V/k}\otimes L).
                                                        \tag{C.3an.9}
\]

We record its integral $K_0$ compatibility in order to
compute both exceptional and local curve classes.
The Zariski descent sequence of
\citet[Theorem~10.3, pp.~383--386]{D_ThomasonTrobaugh1990}
has $E_2^{a,b}=H^a(\mathcal K_b)$ and abuts to $K_{b-a}$.
In dimension at most three the only possible incoming
differential to $H^2(\mathcal K_2)$ comes from
$H^0(\mathcal K_1)$. Every global unit lifts to actual
$K_1$, so this differential is zero. Every possible
outgoing target has cohomological degree at least four.
Thus on regular $Z$ this term is the degree-two quotient
of the rank-and-determinant filtration of $K_0$; on
the two-dimensional $V_2$ it is the kernel of rank and
determinant itself. These identifications are natural
for the restriction used here. On a regular scheme
\[
 \mathrm{CH}^2\simeq F^2K_0/F^3K_0
\]
integrally, with inverse $-c_2$ and multiplier
$(2-1)!=1$; see \citet[Example~15.3.6, p.~298]{frC4Fulton1998}.
Hence no rational Chern character or division by five
enters this comparison.

For $i:V\hookrightarrow Z$ and
$\beta\in\mathrm{CH}^1(V)=\operatorname{Pic}(V)$, choose
the corresponding rank-zero $K_0$ representative on $V$.
The zero-section Koszul formula represents its pushforward
by $\pi^*\beta(1-[L])$. Its restriction to $V_2$ is
pulled back from $V$, and therefore has zero relative
component in (C.3an.8). It follows that $\delta(i_*\beta)=0$.
Chow localization gives a well-defined map on
$\mathrm{CH}^2(X)$.

\emph{4. Ordinary algebraic equivalence and the trace.}
The ordinary subgroup $\mathrm{CH}^2(X)_{\mathrm{alg}}$
is divisible. To verify the precise convention, take an
actual family over a smooth connected parameter curve,
normalize and compactify that curve to $Q$, and close
the horizontal cycles in $X\times Q$. A horizontal
integral cycle is flat over this regular curve. Refined
fibres and proper projection give a homomorphism from
$\mathrm{CH}_0(Q)^0=\operatorname{Pic}^0(Q)(k)$.
The principal-divisor relation follows by pulling the
rational function back to each horizontal component
and pushing its divisor forward. This is exactly the
actual-family convention, also used in C.3o.
The group $\operatorname{Pic}^0(Q)(k)$ is divisible,
including by five, because multiplication on its
abelian variety is surjective on algebraically closed
points. Its images and their finite sums are divisible.
Every homomorphism from that subgroup to a group of
exponent five is consequently zero.

Pair (C.3an.9) with (C.3an.5) and the Serre trace:
\[
 \rho:\mathrm{CH}^2(X)\longrightarrow k,
       \qquad \rho(\xi)=\operatorname{Tr}(\delta\xi\wedge\eta).
                                                        \tag{C.3an.10}
\]
The trace and its local normalization are those of
\citet[Section~48.27, especially Lemma~48.27.1 and
Remarks~48.27.2--3,6, Tag 0FVU]{D_Stacks0FVU}.
The preceding divisibility proves that $\rho$ factors
through the ordinary quotient $B^2(X)$.

\emph{5. Actual curves with every scalar value.}
At the affine point $v:y=z=s=T=0$, choose $s,T$ as
etale parameters and the displayed trivialization of $L$.
Let $t$ be the fibre coordinate on $Z_0$. For $a\in k$
consider
\[
 s=0,\qquad T=at,\qquad z=-a^2t^2,\qquad
                         y^5-y=a^8t^8.
                                                        \tag{C.3an.11}
\]
For $a\ne0$ this affine curve is integral and smooth.
Indeed a rational solution of $h^5-h=a^8t^8$ has no
finite pole, hence is a polynomial; its nonconstant
pole order would be divisible by five, not eight.
The Artin--Schreier polynomial is therefore irreducible.
For $a=0$ take its component $y=0$.
The selected point $t=y=0$ is smooth and transverse
to the zero section because the $y$ derivative is $-1$.

Close the curve in $Z_0$ and take its reduced affine
image $B_a$ in the cone. The proper map of curves is
birational off the vertex and has finite vertex fibre.
Its normalization is the finite normalization of $B_a$.
After the strict-local base change, that finite normal
algebra decomposes into henselian local branches.
Select the branch whose closed point is the specified
$v$. It gives an actual integral curve $D_a\subset Z$.
The normalization is an isomorphism at the selected
smooth germ. Its remaining non-isomorphism locus is
closed and proper over the henselian local branch;
if nonempty its image would contain the closed point,
a contradiction. Thus the selected map is a closed
embedding and meets $V$ only at $v$. This construction
selects an actual branch for each fixed scalar; it
does not assert a simultaneous family in $a$.

The derived intersections with $V_2$ are the perfect
Koszul point classes of $(s,T-at)$ and $(s,T)$.
On the overlap where $sT$ is invertible their difference
has relative symbol and image
\[
 \{s,T-at\}-\{s,T\}\longmapsto
                            t\,\frac{a\,ds}{sT}.
                                                        \tag{C.3an.12}
\]
The two localization boundaries give the Koszul product,
so this is its local-cohomology representative under
the integral descent edge. Wedging with
$\eta=t^{-1}dT$ gives the local residue $a$.
Functorial Serre duality normalizes this trace: the
Koszul generator of
$\operatorname{Ext}^2(\mathcal O_v,\Omega^2_V)$ maps
to $H^2(\Omega^2_V)$ with trace dual to the identity
$H^0(\mathcal O_V)\to H^0(\mathcal O_v)=k$.
In the ordered parameters $s,T$ its representative is
$ds\wedge dT/(sT)$, with trace one. Consequently
\[
                         \rho([D_a]-[D_0])=a.
                                                        \tag{C.3an.13}
\]
Every scalar is attained, proving the required surjection.
\end{proof}

\emph{Source-sign audit and scope.}
The opened RIMS1736 preprint prints a plus sign in its
equation~(7) and repeats it in the cyclic-cover root on
author p.~7. With its extension and line-section conventions
this conflicts with the positive normal bundle and its own
author p.~14 formula $G\sim p(S_\infty-D_1)$.
RIMS1505, author pp.~2--4, explicitly derives
$\mathcal O_P(pS_\infty-G)=\bar g^*\mathcal O_C(pD_1)$,
although its subsequent cover formula also has a sign
inconsistency. Equations (C.3an.2)--(C.3an.3) were therefore
proved independently by fibre degree, restriction and the
typed root. The published 2013 article is identified
bibliographically, but its full publisher PDF was blocked;
no claim about its printed signs is made.
Only infinitesimal restriction of actual curves is used.
There is no algebraization of arbitrary formal cycles,
rational Chern character, or positive-characteristic
Hodge or Kodaira vanishing assumption.
